\begin{textsample}{POS Dim 1 – 2013\_c5 – Score 10.00 – 2013\_c5/enade\_2013\_odonto\_ed444be374.txt}  \label{ex:f1_pos_083}
CAPÍTULO 5 DISTRIBUIÇÃO DOS CONCEITOS 5.1 \textbf{PANORAMA} NACIONAL DA DISTRIBUIÇÃO DOS CONCEITOS A Tabela 5.1 apresenta a quantidade e distribuição de cursos de Odontologia participantes do ENADE/2013, por faixa de conceito e Grande Região.
A diferença entre os cursos \textbf{tabulados} neste capítulo e no capítulo 2 são os cursos sem conceito, em princípio, aqueles sem alunos concluintes que participassem da prova ou que tivessem tão somente um aluno concluinte13.
Observando -se os dados da Tabela 5.1, nota -se que, dos 184 cursos participantes, 71 ( 38,6 \% ) classificaram -se com conceito 3, o valor modal.
Este foi também o conceito modal em todas as regiões do Brasil ( entre 36,8 \% na região Norte e 44,7 \% na região Sudeste ), exceto na região Nordeste, onde a moda foi o conceito 4 ( 27,3 \% ).
O conceito 2 foi o segundo mais frequente em nível nacional ( 26,1 \%, correspondendo a 48 cursos ) e o conceito 4, o terceiro ( 19,0 \%, correspondendo a 35 cursos ).
Houve, ainda, 16 cursos ( 8,7 \% ) que receberam conceito 5 e outros nove cursos que receberam conceito 1.
Dos 184 cursos de Odontologia, cinco ficaram sem conceito ( SC ).
A região Norte participou com 19 cursos ou 10,3 \% do total nacional.
Nesta região, o conceito 3 ( conceito modal ) foi atribuído a sete cursos dos 19 participantes, o que \textbf{equivale} a 36,8 \% do total regional, como já comentado.
Apenas um curso ( 5,3 \% ) ficou sem conceito(SC ).
Os demais foram avaliados com conceito 1 ( três cursos, correspondendo a 15,8 \% ), conceito 2 ( seis cursos correspondendo a 31,6 \% ) e conceito 4 ( dois cursos correspondendo a 10,5 \% ).
Nenhum curso foi alocado ao conceito 5. 13 Estes cursos com somente um concluinte não foram considerados no capítulo 2.
A região Nordeste participou com 33 cursos ou 17,9 \% do total nacional.
Nesta região, nove cursos, 27,3 \% em termos regionais, obtiveram conceito 4, o conceito modal para a região.
Os conceitos 2 e 3 foram atribuídos a, respectivamente, oito e sete cursos ( 24,2 \% e 21,2 \% ).
O conceito 5 foi atribuído a, também, oito cursos.
Nessa região, apenas um dos cursos ( 3,0 \% ) ficou sem conceito ( SC ), enquanto nenhum recebeu conceito 1.
Como já comentado, dos 85 cursos participantes da região Sudeste, 38 ( 44,7 \% ) obtiveram conceito 3, o conceito modal.
O conceito 1 foi atribuído a quatro cursos ( 4,7 \% ) e o conceito 2, a 25 ( 29,4 \% ).
O conceito 4 foi atribuído a 11 cursos ( 12,9 \% ) e quatro outros ( 4,7 \% ) receberam o conceito 5.
Apenas três cursos ( 3,5 \% ) ficaram sem conceito ( SC ).
A região Sul também \textbf{contou} com cursos distribuídos em todas as faixas de conceitos, exceto no conceito 1.
A predominância do conceito 3 foi de 39,4 \%, como já comentado anteriormente, correspondentes a 13 dos 33 cursos participantes na região Sul.
O conceito 2 foi atribuído a oito cursos ( 24,2 \% ) e o conceito 4, a dez cursos ( 30,3 \% ).
Nenhum curso ficou sem conceito ( SC ) ou recebeu conceito 1, enquanto o conceito 5, foi atribuído a dois deles ( 6,1 \% ).
Como já comentado, de todos os 14 cursos participantes na região Centro-Oeste, seis receberam conceito 3, o conceito modal ( 42,9 \% ).
Nesta região nenhum curso ficou sem conceito ( SC ).
Já, dois cursos foram alocados no conceito 1 e outros dois no conceito 5, correspondendo a 28,6 \% dos cursos da região ( 14,3 \%, cada curso ).
Os demais cursos foram avaliados com conceito 2 ( um curso, 7,1 \% ) e conceito 4 ( três cursos, 21,4 \% ).
O Gráfico 5.1 apresenta a distribuição do Conceito ENADE segundo Grande Região : \textbf{poligonais} mais à esquerda representam Grandes Regiões com pior distribuição, e \textbf{poligonais} mais à direita, Grandes Regiões com melhores conceitos.
As regiões Norte ( linha azul ) e Sudeste ( linha vermelha ) apresentaram os cursos em IES com valores mais baixos, ao passo que a regiões Nordeste ( linha verde ), Sul ( linha \textbf{roxa} ) e Centro-Oeste ( linha preta ) apresentaram as distribuições correspondendo a cursos com valores maiores. 5.2 CONCEITOS POR CATEGORIA ADMINISTRATIVA E POR GRANDE REGIÃO A Tabela 5.2 apresenta a distribuição dos cursos participantes do ENADE/2013 de Odontologia, por Categoria Administrativa, de acordo com os conceitos por eles alcançados, segundo as Grandes Regiões brasileiras.
Dos 184 cursos participantes, 54 ( 29,3 \% ) eram \textbf{ministrados} em Instituições \textbf{Públicas} e 130 ( 70,7 \% ), em privadas.
De acordo com as informações da Tabela 5.2, em termos nacionais, das 16 IES que receberam o conceito 5, 11 eram públicas e cinco privadas.
Dos 54 cursos participantes de IES Públicas, o conceito 4 foi o valor modal, atribuído a 28 cursos ( 51,9 \% ).
Entre os demais cursos de Instituições \textbf{Públicas} participantes, um recebeu conceito 1 ( 1,9 \% da categoria ), três obtiveram conceito 2 ( 5,6 \% ) e 22 foram igualmente alocados aos conceitos 3 e 5 ( 20,4 \%, cada ).
Nesta categoria, nenhum dos cursos ficou sem conceito ( SC ).
Na rede privada, o conceito modal foi 3, com 60 cursos dos 130 da categoria.
Entre os demais cursos participantes, oito receberam conceito 1 e 45, conceito 2.
O conceito 4 foi atribuído a atribuído a sete cursos e o conceito 5, como já comentado, a 11 cursos.
Nesta Categoria Administrativa, cinco dos cursos ficaram sem conceito ( SC ).
Na análise por região, observa -se que, na região Norte, as Instituições Privadas participaram com 15 cursos ( 78,9 \% do total regional, foi a mais elevada do que nas demais regiões brasileiras ), dos quais somente um curso ficou sem conceito ( SC ).
O conceito modal para as Instituições Privadas na região foi 2, com seis cursos, correspondendo a 40,0 \%.
O conceito 3 apresentou cinco cursos, correspondendo a 33,3 \%.
Outros dois cursos receberam conceito 1 ( 13,3 \% ).
Somente um curso foi alocado no conceito 4 ( 6,7 \% ), enquanto nenhum foi alocado no conceito 5.
As Instituições Públicas participaram com quatro cursos ( 21,1 \% do total regional ), dos quais apenas um ficou com conceito 1.
O conceito modal para esta Categoria Administrativa, na região Norte, foi 3, com dois cursos, correspondendo a 50,0 \%.
Os conceitos 1 e 4 participaram, também, com um curso cada ( 25,0 \%, cada ).
Nenhum curso ficou sem conceito ( SC ) ou foi alocado aos conceitos 2 ou 5 nesta categoria e região.
Na região Nordeste, a rede privada concentrou 16 dos 33 cursos participantes, equivalentes a 48,5 \% do total da região, a menor proporção entre todas as regiões.
Nesta Categoria Administrativa na região, predominaram os cursos com conceito 2, oito cursos correspondendo a 50,0 \% dos cursos oferecidos por IES Privadas no Nordeste.
Nenhum curso recebeu o conceito 1 ou recebeu o conceito 5.
Os demais foram avaliados com conceitos 3 ( seis cursos correspondendo a 37,5 \% ) e 4 ( um curso correspondendo a 6,3 \% ).
Nessa categoria e região, somente um curso ficou sem conceito ( SC ).
As Instituições \textbf{Públicas} dessa região participaram com 17 cursos ( 51,5 \% ), dos quais 16 obtiveram conceitos 4 e 5 ( oito cursos, cada ), os valores modais ( \textbf{bimodal} ).
Já o conceito 3 foi alocado a somente um curso, correspondendo a 5,9 \%.
Nenhum curso ficou sem conceito ( SC ) ou recebeu os conceitos 1 ou 2 nesta combinação de categoria e região.
Na região Sudeste, a proporção de cursos da rede privada, 78,8 \%, correspondendo a 67 dos 85 cursos participantes.
Nesta categoria, na região Sudeste, o conceito modal foi 3 ( 32 cursos ).
Nesta combinação de Categoria Administrativa e Grande Região, apenas três dos cursos não receberam conceito ( SC ).
Os demais foram avaliados com conceito 1 ( quatro cursos ), conceito 2 ( 22 cursos ), conceito 4 ( dois cursos ) e conceito 5 ( também, quatro cursos ).
Entre os 18 cursos em Instituições \textbf{Públicas} na região Sudeste, a categoria modal foi a que recebeu conceito 4 ( nove cursos ).
Os demais foram avaliados com os conceitos 2 ( três cursos ) e 3 ( seis cursos ).
Na região Sudeste, nesta categoria, nenhum curso ficou sem conceito ( SC ) ou recebeu os conceitos 1 ou 5.
As Instituições Privadas concentraram 21 dos 33 cursos participantes da região Sul, 63,6 \% do total regional.
Desses, 11 ficaram com conceito 3, o conceito modal.
Nesta combinação de Categoria Administrativa e Grande Região, nenhum curso ficou sem conceito ( SC ) ou recebeu os conceitos 1 ou 5, enquanto oito cursos receberam conceito 2.
Apenas dois cursos foram alocados ao conceito 4.
As Instituições \textbf{Públicas} na região Sul participaram com 12 cursos ( 36,4 \% ), dos quais nenhum ficou sem conceito ( SC ) ou recebeu os conceitos 1 ou 2.
O conceito modal foi 4, com oito cursos.
Quatro cursos foram alocados igualmente aos conceitos 3 e 5 ( dois cursos, cada ).
Na região Centro-Oeste, 11 dos 14 cursos participantes eram de Instituições Privadas ( 78,6 \% em termos regionais ).
Destes, seis concentraram -se no conceito 3, conceito modal.
Os demais receberam conceito 1 ( dois cursos ) e conceitos 2, 4 e 5 ( um curso, cada ).
Nenhum curso ficou sem conceito ( SC ) nesta categoria.
Dos três cursos de Instituições \textbf{Públicas}, dois foram avaliados no conceito 4, conceito modal, e um único curso foi alocado ao conceito 5.
Nesta região, nenhum curso de IES \textbf{pública} ficou sem conceito ( SC ) ou recebeu os demais conceitos \textbf{restantes} nessa categoria.
O Gráfico 5.2 apresenta a distribuição do Conceito ENADE segundo a Categoria Administrativa da IES.
Os cursos em IES Públicas ( linha azul ) apresentam uma \textbf{poligonal} mais à direita do que aqueles em IES Privadas ( linha verde ) e consequentemente uma distribuição de Conceitos ENADE com valores maiores. 5.3 CONCEITOS POR ORGANIZAÇÃO ACADÊMICA E POR GRANDE REGIÃO Na Tabela 5.3 encontra -se a distribuição dos conceitos atribuídos aos cursos participantes do ENADE/2013 na Área de Odontologia, por Organização Acadêmica, segundo as Grandes Regiões brasileiras.
Dos 184 cursos participantes, 121 eram oferecidos em Universidades, 41 em Faculdades e os demais 22 em Centros Universitários.
Esta distribuição corresponde a, respectivamente, 65,8 \%, 22,3 \% e 12,0 \% dos cursos.
De acordo com os dados apresentados, de todos os 16 cursos avaliados com conceito 5, 14 eram vinculados a Universidades.
Somente quatro cursos em Universidades ficaram sem conceito ( SC ).
Esse tipo de Organização Acadêmica teve o conceito 3 como modal, com 43 cursos ( 35,5 \% ).
Os demais cursos de Universidades avaliados, receberam os conceitos 1 ( dois cursos ), 2 ( 25 cursos ), 4 ( 33 cursos ) e conceito 5 ( 14 cursos, como já mencionado ).
Entre os cursos em Centros Universitários, o conceito modal também foi 3, com dez cursos ( 45,5 \% ).
Neste tipo de organização acadêmica nenhum curso ficou sem conceito ( SC ) ou recebeu os conceitos 4 ou 5.
Os outros cursos neste tipo de Organização Acadêmica receberam os conceitos 1 ( três cursos ) e 2 ( nove cursos ).
Nas Faculdades, somente um dos 41 cursos ficou sem conceito ( SC ) e 18 ( 43,9 \% ) receberam o mesmo conceito modal dos demais tipos de Organização Acadêmica, conceito 3.
Dos demais cursos neste tipo de Organização Acadêmica, quatro receberam conceito 1, 14 conceito 2, e quatro, conceitos 4 e 5 ( dois cursos, cada ).
Considerando -se separadamente as regiões brasileiras, verifica -se que, na região Norte, as Universidades concentraram cinco dos 19 cursos participantes.
Dos cursos em Universidades, nenhum ficou sem conceito ( SC ) ou recebeu conceitos 1 ou 5.
Quatro cursos ficaram alocados igualmente aos conceitos 3 e 4 ( dois cursos, cada ), os conceitos modais, enquanto somente um curso recebeu conceito 2 nessa categoria e região.
Os Centros Universitários da região Norte foram representados por apenas três cursos, destes, todos ficaram alocados igualmente aos conceitos 1, 2 e 3 ( um curso, cada ).
Nenhum curso ficou sem conceito ( SC ) ou recebeu os conceitos 4 ou 5.
As Faculdades participaram com 11 cursos na região Norte.
Nenhum curso ficou com os conceitos 4 ou 5.
Os demais receberam conceito 1 ( dois cursos ) e conceitos 2 e 3 ( quatro cursos cada, conceito modal ).
Somente um curso ficou sem conceito ( SC ).
Na região Nordeste, as Universidades participaram com 21 dos 33 cursos na Área de Odontologia da região.
Nenhum dos cursos em Universidades no Nordeste ficou com conceito 1.
Nessa região, 16 cursos foram alocados igualmente aos conceitos 4 e 5 ( oito cursos, cada ), conceitos modais ( \textbf{bimodal} ).
Os demais, receberam conceito 2 ( um curso ) e 3 ( três cursos ).
Somente um curso ficou sem conceito ( SC ) nessa categoria e região.
Os Centros Universitários \textbf{contaram} com três cursos participantes na região Nordeste, todos eles concentrados no conceito 2 representado a totalidade da categoria.
As Faculdades foram representadas por nove cursos na região Nordeste, dos quais oito cursos foram alocados igualmente aos conceitos 2 e 3 ( quatro cursos, cada ), os conceitos modais ( \textbf{bimodal} ) na categoria.
Somente um curso recebeu o conceito 4.
Nenhum curso ficou sem conceito ( SC ) ou recebeu os conceitos 1 ou 5.
Na região Sudeste, as Universidades concentraram 60 dos 85 cursos da região.
Entre os cursos em Universidades na região, o conceito modal foi 3 com 24 cursos, e apenas três cursos ficaram sem conceito ( SC ), enquanto dois cursos receberam conceito 1.
Os demais cursos receberam os conceitos 2 ( 18 cursos ), 4 ( dez cursos ) e 5 ( três cursos ).
Os Centros Universitários participaram com 11 cursos na região Sudeste, dos quais seis obtiveram, também, o conceito modal 3.
Nenhum ficou sem conceito ( SC ) ou recebeu os conceitos 4 ou 5.
Somente um curso foi alocado ao conceito 1.
As Faculdades foram representadas por 14 cursos na região Sudeste, que se distribuíram nos conceitos 1, 4 e 5 ( um curso, cada ), 2 ( três cursos ) e 3 ( oito cursos, conceito modal ).
Nenhum curso ficou sem conceito ( SC ) nessa categoria e região.
Dos 33 cursos da região Sul, 27 eram de Universidades, para os quais os conceitos modais foram 3 e 4 ( \textbf{bimodal} ), com dez cursos, cada.
Nesse tipo de organização, nenhum dos cursos ficou sem conceito ( SC ) ou recebeu conceito 1.
Os demais receberam os conceitos 2 ( cinco cursos ) e 5 ( dois cursos ).
Os Centros Universitários da região Sul tiveram dois cursos participantes todos alocados ao conceito 3, configurando, portanto, a totalidade na categoria.
Foram quatro os cursos vinculados a Faculdades na região Sul, dos quais, três receberam conceito 2, o conceito modal, enquanto somente um curso ficou alocado ao conceito 3.
Nenhum curso obteve conceito ( SC ) nessa categoria e região.
Na região Centro-Oeste, oito dos 14 cursos eram de Universidades.
Nesse tipo de organização, nenhum curso ficou sem conceito ( SC ) ou recebeu os conceitos 1 ou 2.
Já o conceito modal foi 3, com quatro cursos.
Os outros cursos obtiveram os conceitos 4 ( três cursos ) e 5 ( um curso ).
Os Centros Universitários da região Centro-Oeste \textbf{contaram} com três cursos, todos alocados igualmente aos conceitos 1, 2 e 3 ( um curso, cada ).
Nenhum curso ficou sem conceito ( SC ) ou recebeu os conceitos 4 ou 5.
Dos três cursos em Faculdades na região Centro-Oeste, todos, também, alocados igualmente, porém, aos conceitos 1, 3 e 5 ( um curso, cada ).
Nessa categoria e região, nenhum curso ficou sem conceito ( SC ) ou recebeu os conceitos 2 ou 4.
O Gráfico 5.3 apresenta a distribuição do Conceito ENADE segundo a Organização Acadêmica da IES.
Os cursos em Universidades ( linha azul ) apresentam uma \textbf{poligonal} mais à direita do que as demais, \textbf{denotando} uma melhor distribuição de Conceitos ENADE.
No outro extremo temos os cursos em Centros Universitários ( linha verde ) com a \textbf{poligonal} mais à esquerda e consequentemente a distribuição de Conceitos ENADE com menores valores.
Cursos em Faculdades ( linha vermelha ) apresentam uma situação intermediária, porém com maior similaridade com aqueles em Centro Universitários ( linha verde ).

% matched lemmas: bimodal, contar, denotar, equivaler, ministrar, panorama, poligonal, pública, restante, roxa, tabular
\end{textsample}
